\documentclass[11pt]{article}

\usepackage[utf8]{inputenc}
\usepackage[T1]{fontenc}
\usepackage{lmodern}
\usepackage[margin=1in]{geometry}
\usepackage{microtype}
\usepackage{enumitem}
\usepackage{titlesec}
\usepackage{xcolor}

\definecolor{accent}{HTML}{1A5276}   % deep blue for headings
\definecolor{linkblue}{HTML}{1A5FB4}  % clickable links
\definecolor{notegray}{HTML}{454545}  % annotation text

\usepackage{hyperref}
\hypersetup{
  colorlinks=true,
  linkcolor=linkblue,
  urlcolor=linkblue,
  citecolor=linkblue,
  pdftitle={PKU-Zurich PhD Summer School 2026 -- Reading List},
  pdfauthor={PKU-Zurich PhD Summer School}
}

% --- Section styling -------------------------------------------------------
\titleformat{\section}{\large\bfseries\color{accent}}{}{0em}{}
\titlespacing*{\section}{0pt}{1.5em}{0.5em}

\setlist[itemize]{leftmargin=1.4em, itemsep=7pt, topsep=4pt}

% --- A reading entry: {authors (year)}{title}{venue}{url}{note} ------------
\newcommand{\reading}[5]{%
  \item #1. \textit{#2}. #3.\ \href{#4}{[link]}%
}

\newcommand{\bglabel}{\smallskip\noindent\textbf{\color{accent}Recommended background}\par\nopagebreak\smallskip}

\pagestyle{plain}

\begin{document}

% ===========================================================================
\begin{center}
  {\LARGE\bfseries\color{accent} PKU--Zurich PhD Summer School}\\[3pt]
  {\large\bfseries on Machine Learning for Macroeconomics and Finance}\\[10pt]
  {\Large\bfseries Reading List}\\[8pt]
  {\normalsize Beijing, China \quad$\cdot$\quad July 6--10, 2026}
\end{center}

\vspace{0.4em}
{\color{accent}\hrule height 1pt}
\vspace{1.2em}

\noindent

% ===========================================================================
\section*{Before the Summer School (Pre-Reading \& Pre-Coding)}

Everyone should complete these two items before arriving in Beijing.

\begin{itemize}
  \item \textbf{Reading --- heterogeneous-agent models.}
  Dirk Krueger, \textit{Heterogeneous Agent Macroeconomics} (lecture notes), 2025.
  \href{https://www.uni-bielefeld.de/fakultaeten/wirtschaftswissenschaften/einrichtungen/bigsem/profiles/economics/winter-term-2025-2026/HeteroBookLatex2025.pdf}{[link]}\\[1.5pt]
  {\small\color{notegray}Read \textbf{Chapter 6, ``The SIM in General Equilibrium,''} carefully; use \textbf{Chapter 4, ``The SIM in Partial Equilibrium,''} as a background reference.}

  \item \textbf{Coding --- neural networks for economic models.}
  Simon Scheidegger, \textit{Deep Equilibrium Nets: Brock--Mirman tutorials} (Jupyter notebooks).\\[1.5pt]
  {\small\color{notegray}Read and run both notebooks:
  Notebook 1 ---
  \href{https://github.com/sischei/Deep_Learning_for_Solving_And_Estimating_Dynamic_Economic_Models/blob/main/lectures/lecture_03_deep_equilibrium_nets/code/lecture_03_01_Brock_Mirman_1972_DEQN.ipynb}{deterministic Brock--Mirman},
  and Notebook 2 ---
  \href{https://github.com/sischei/Deep_Learning_for_Solving_And_Estimating_Dynamic_Economic_Models/blob/main/lectures/lecture_03_deep_equilibrium_nets/code/lecture_03_02_Brock_Mirman_Uncertainty_DEQN.ipynb}{stochastic Brock--Mirman}.}
\end{itemize}

% ===========================================================================
\section*{Day 1 --- Deep Learning and Heterogeneous-Agent Basics}
\textit{Serguei Maliar (Santa Clara University)}
\smallskip

\begin{itemize}
  \reading{Maliar, Maliar \& Winant (2021)}{Deep learning for solving dynamic economic models}{Journal of Monetary Economics 122, 76--101}{https://doi.org/10.1016/j.jmoneco.2021.07.004}{}
  \reading{Lepetyuk, Maliar \& Maliar (2020)}{When the U.S. catches a cold, Canada sneezes: A lower-bound tale told by deep learning}{Journal of Economic Dynamics and Control 117, 103926}{https://ideas.repec.org/a/eee/dyncon/v117y2020ics0165188920300944.html}{}
  \reading{Maliar \& Maliar (2021)}{Deep learning classification: Modeling discrete labor choice}{Journal of Economic Dynamics and Control 135, 104295}{https://ideas.repec.org/a/eee/dyncon/v135y2022ics016518892100230x.html}{}
\end{itemize}

\bglabel
\begin{itemize}
  \reading{Goodfellow, Bengio \& Courville (2016)}{Deep Learning}{MIT Press (full text free online)}{https://www.deeplearningbook.org/}{The canonical deep-learning textbook. Skim Part~I and Chapters 6--8 (feedforward networks, regularization, optimization/SGD).}
  \reading{Aiyagari (1994)}{Uninsured Idiosyncratic Risk and Aggregate Saving}{Quarterly Journal of Economics 109(3), 659--684}{https://academic.oup.com/qje/article-abstract/109/3/659/1838287}{The foundational discrete-time, incomplete-markets heterogeneous-agent model behind the discrete-time HA lecture.}
  \reading{Krusell \& Smith (1998)}{Income and Wealth Heterogeneity in the Macroeconomy}{Journal of Political Economy 106(5), 867--896}{https://doi.org/10.1086/250034}{The classic heterogeneous-agent model with aggregate risk; its curse of dimensionality is exactly what the deep-learning solvers are designed to overcome.}
\end{itemize}

% ===========================================================================
\section*{Day 2 --- Deep Equilibrium Nets}
\textit{Simon Scheidegger (HEC Lausanne) and Felix Kübler (University of Zurich)}
\smallskip

\begin{itemize}
  \reading{Scheidegger (2026)}{Deep Learning for Solving and Estimating Dynamic Models in Economics and Finance}{arXiv:2605.14493 (lecture notes, with TensorFlow / PyTorch notebooks)}{https://arxiv.org/abs/2605.14493}{}
  \reading{Azinovic, Gaegauf \& Scheidegger (2022)}{Deep Equilibrium Nets}{International Economic Review 63(4), 1471--1525. Code: \href{https://github.com/sischei/DeepEquilibriumNets}{github.com/\allowbreak sischei/\allowbreak DeepEquilibriumNets}}{https://doi.org/10.1111/iere.12575}{The flagship paper the DEQN lectures are built on: neural networks trained unsupervised to satisfy equilibrium conditions along simulated paths, with heterogeneity, uncertainty, and occasionally binding constraints. Reproduction code at \href{https://github.com/sischei/DeepEquilibriumNets}{github.com/sischei/DeepEquilibriumNets}.}
  \reading{Scheidegger \& Bilionis (2019)}{Machine learning for high-dimensional dynamic stochastic economies}{Journal of Computational Science 33, 68--82}{https://doi.org/10.1016/j.jocs.2019.03.004}{Combines Gaussian-process regression with active subspaces to solve models up to roughly 500 dimensions; bridges the deep-learning and Gaussian-process halves of Day 2.}
  \reading{Chen, Didisheim \& Scheidegger (2026)}{Deep Surrogates for Finance: With an Application to Option Pricing}{Journal of Financial Economics 177, 104222}{https://doi.org/10.1016/j.jfineco.2025.104222}{Underpins the ``deep surrogate models and uncertainty quantification'' lecture: neural-network surrogates approximate expensive structural models to accelerate estimation by orders of magnitude.}
  \reading{Friedl, Kübler, Scheidegger \& Usui (2021)}{Deep Uncertainty Quantification: With an Application to Integrated Assessment Models}{Working paper, University of Lausanne}{https://sites.google.com/site/simonscheidegger/working-papers}{The ``deep uncertainty quantification'' counterpart to the equilibrium-net methods: a deep-learning global solver paired with a Gaussian-process surrogate to perform parametric uncertainty quantification in stochastic dynamic models, illustrated on climate integrated-assessment models.}
  \reading{Rasmussen \& Williams (2006)}{Gaussian Processes for Machine Learning}{MIT Press (full text free online)}{https://gaussianprocess.org/gpml/}{The canonical reference for Kübler's Gaussian-processes and Bayesian numerical-methods lecture; focus on Chapters 2--5 (GP regression, covariance functions, hyperparameter learning).}
\end{itemize}

% ===========================================================================
\section*{Day 3 --- DeepHAM and Structural Reinforcement Learning}
\textit{Yucheng Yang (University of Zurich), Benjamin Moll (LSE)}
\smallskip

\begin{itemize}
  \reading{Han, Yang \& E (2026)}{DeepHAM: A Global Solution Method for Heterogeneous Agent Models with Aggregate Shocks}{Quantitative Economics 17(2), 297--341}{https://doi.org/10.3982/QE2190}{The DeepHAM method: neural networks approximate value and policy functions and represent the cross-sectional distribution via learned generalized moments. Focus on the moment representation and the simulation-based training objective.}
  \reading{Yang, Wang, Schaab \& Moll (2025)}{Structural Reinforcement Learning for Heterogeneous Agent Macroeconomics}{arXiv:2512.18892 (code coming soon)}{https://arxiv.org/abs/2512.18892}{The basis for Moll's special online lecture: a structural RL method that uses low-dimensional prices as state variables and lets agents learn equilibrium price dynamics from simulated paths (Krusell--Smith, Huggett with aggregate shocks, HANK). Code coming soon.}
  \reading{Wibault et al. (2026)}{Recurrent Structural Policy Gradient for Partially Observable Mean Field Games}{International Conference on Machine Learning (ICML), 2026. Code: \href{https://github.com/CWibault/mfax}{github.com/\allowbreak CWibault/\allowbreak mfax}}{https://arxiv.org/abs/2602.20141}{A history-aware ``structural'' policy-gradient method for partially observable mean-field games with common noise that solves, for the first time, a macroeconomics mean-field game with heterogeneous agents and history-dependent policies; introduces the MFAX (JAX) framework. Code at \href{https://github.com/CWibault/mfax}{github.com/CWibault/mfax}.}
\end{itemize}

\bglabel
\begin{itemize}
  \reading{Moll (2025)}{The Trouble with Rational Expectations in Heterogeneous Agent Models: A Challenge for Macroeconomics}{The Economic Journal (forthcoming); Economic Journal Lecture, RES 2024}{https://arxiv.org/abs/2508.20571}{Argues that rational expectations about equilibrium prices --- which force agents to forecast the entire cross-sectional distribution (the ``master equation'') --- are computationally and behaviorally implausible, and points to learning and reinforcement-learning alternatives; the conceptual motivation for the structural-RL lecture.}
\end{itemize}

% ===========================================================================
\section*{Day 4 --- Deep Learning and Continuous-Time Macro-Finance}
\textit{Goutham Gopalakrishna (University of Toronto, Rotman)}
\smallskip

\begin{itemize}
  \reading{Gopalakrishna \& Wu (2021)}{ALIENs for Continuous Time Economies}{Swiss Finance Institute Research Paper No. 21-34}{https://ssrn.com/abstract=3848657}{The instructor's flagship method: Active-Learning-Inspired Equilibrium Nets solve high-dimensional, nonlinear continuous-time macro-finance models by combining time-stepping with active learning over the economically important regions of the state space.}
  \reading{Wu, Guo, Gopalakrishna \& Poulos (2024)}{Deep-MacroFin: Informed Equilibrium Neural Network for Continuous Time Economic Models}{arXiv:2408.10368. Code: \href{https://github.com/rotmanfinhub/deep-macrofin}{github.com/\allowbreak rotmanfinhub/\allowbreak deep-macrofin}}{https://arxiv.org/abs/2408.10368}{The companion library and method paper; the natural reference for the hands-on tutorial, showing how to encode HJB and equilibrium conditions as neural-network losses. Code at \href{https://github.com/rotmanfinhub/deep-macrofin}{github.com/rotmanfinhub/deep-macrofin}.}
  \reading{Sirignano \& Spiliopoulos (2018)}{DGM: A deep learning algorithm for solving partial differential equations}{Journal of Computational Physics 375, 1339--1364}{https://doi.org/10.1016/j.jcp.2018.08.029}{The Deep Galerkin Method behind the ``deep learning for PDEs'' lecture: a meshfree algorithm that trains a neural network to satisfy a high-dimensional PDE on randomly sampled points.}
\end{itemize}

\bglabel
\begin{itemize}
  \reading{Brunnermeier \& Sannikov (2014)}{A Macroeconomic Model with a Financial Sector}{American Economic Review 104(2), 379--421}{https://www.aeaweb.org/articles?id=10.1257/aer.104.2.379}{The foundational continuous-time macro-finance model with a financial sector; its full nonlinear crisis dynamics are the canonical problem class these solvers target.}
  \reading{Achdou, Han, Lasry, Lions \& Moll (2022)}{Income and Wealth Distribution in Macroeconomics: A Continuous-Time Approach}{Review of Economic Studies 89(1), 45--86}{https://academic.oup.com/restud/article/89/1/45/6149490}{Recasts the Aiyagari--Bewley--Huggett model in continuous time as coupled HJB and Kolmogorov-forward PDEs --- the continuous-time HA toolkit underlying the lectures.}
\end{itemize}

% ===========================================================================
\section*{Day 5 --- Deep Learning for Macro-Finance and Beyond}
\textit{Jonathan Payne (Princeton University); Keynote: Weinan E (Peking University)}
\smallskip

\begin{itemize}
  \reading{Gu, Laurière, Merkel \& Payne (2024)}{Global Solutions to Master Equations for Continuous Time Heterogeneous Agent Macroeconomic Models}{arXiv:2406.13726 (forthcoming, JPE: Macroeconomics)}{https://arxiv.org/abs/2406.13726}{Introduces Economic-Model-Informed Neural Networks (EMINNs) to compute global solutions of master equations, with the continuous-time Krusell--Smith model as the lead application.}
  \reading{Payne, Rebei \& Yang (2025)}{Deep Learning for Search and Matching Models}{Swiss Finance Institute Research Paper No. 25-05}{https://ideas.repec.org/p/chf/rpseri/rp2505.html}{Underpins the search-and-matching lecture: extends the deep-learning-for-PDE approach to labor-market models with aggregate shocks and heterogeneous agents.}
  \reading{Gopalakrishna, Gu \& Payne (2026)}{Institutional Asset Pricing with Segmentation and Household Heterogeneity}{SSRN Working Paper}{https://ssrn.com/abstract=6162106}{The basis for the asset-pricing-with-household-heterogeneity lecture: a segmented heterogeneous-agent macro-finance model solved globally by deep learning, calibrated to Flow of Funds data.}
\end{itemize}

\vfill
\begin{center}
{\footnotesize\color{notegray}
Summer School website: \href{https://liboecon.com/2026_summer_school.html}{liboecon.com/2026\_summer\_school.html}
\quad$\cdot$\quad Last updated: \today}
\end{center}

\end{document}
